\documentclass[letterpaper,11pt,leqno]{article}
\usepackage{paper}
\bibliographystyle{paper}

% Enter paper title to populate PDF metadata:
\hypersetup{pdftitle={Minimalist LaTeX Template for Academic Papers}}

% Enter path to BibTeX file with references:
\newcommand{\bib}{paper.bib}

% Enter path to PDF file with figures:
\newcommand{\pdf}{figures.pdf}

\begin{document}

% Enter title:
\title{}

% Enter authors:
\author{Edward Melville
%
% Enter affiliations and acknowledgements:
\thanks{Edward Melville: Westcliff High School for Boys. Email: \href{mailto:melvilleted@hotmail.com}}}

% Enter date:
\date{October 2026}   

% Enter paper URL (can be commented out):
%\paperurl{https://github.com/pmichaillat/latex-paper}

\begin{titlepage}
\maketitle

% Enter abstract:
Low Earth orbit (LEO) has become significantly more congested over recent decades. This is down to a large increase in the number of satellites launched into LEO, but also the detection of smaller fragments. As a result, traditional conjunction assessment has fallen behind computationally. Many organisations are looking for machine learning (ML) to step in. This paper investigates the effectiveness of an ML model as a first stage filter for conjunction screening. Trained on real-world historical conjunction data, the distinction that this paper explores is the model’s ability to distinguish between satellite pairs that should generate a conjunction data message (CDM) and those pairs that occupy broadly similar orbital regions, despite not generating a CDM when propagated. To achieve this, 2 datasets were produced (containing only a satellite pairs’ orbital features). The positive dataset contained satellite pairs that had historically flagged a CDM. The negative dataset contained satellite pairs that had similar orbital features but did not generate a CDM within 7-days. The ML model, which was built on a random forest framework, was then evaluated on one key metric: recall rate. The recall rate of the ML model was determined by the number of incorrectly predicted ‘safe’ satellites – chosen specifically in order to operate a safe-heavy system, as a false positive was far less dangerous than a false negative in our context. When faced with a generic set of just under 300 satellite pairs, the model’s recall rate was 99.65\% resulting in 1 false negative encounter. When the dataset was then altered to match stricter non-conjunction criteria, the model’s recall rate fell to 97.80\% with 6 false negative predictions. Whilst this shows promise that an ML model can intervene as a first-stage filter in the conjunction assessment framework, it also suggests that further development of neural networks in the satellite safety context is required. Understandably so as 1 false negative prediction could cost a satellite operator millions of dollars. 

\end{titlepage}

% Enter main text:
\section{Introduction}\label{s:introduction}
 
Over the past two decades, the number of tracked space objects has almost tripled [1]. It is therefore unsurprising to find that low Earth orbit (LEO) has thus become significantly more congested. The operational defence (conjunction screening and collision avoidance manoeuvres) against catastrophic impacts—which risk triggering an irreversible cascade of orbit debris known as the Kessler Syndrome—is now under a significant computational burden [2]. As the catalogue of orbital objects expands so much that traditional mathematical based propagation can no longer keep up—both from an increased number of space objects launched into space each year but also the detection of smaller objects—the need for a new computational method is required [3]. Many organisations that deal with collision avoidance (COLA) such as Advanced Space and GMV, are now turning to machine learning (ML) to solve these problems. This paper explores the efficiency of these ML models in identifying potentially dangerous satellite pairs for further high-fidelity screening just from their orbital characteristics alone, shifting the COLA timeframe further in advance of any potential incidents.
The advent of a new space era, pioneered by mega constellations such as SpaceX’s Starlink, has permanently altered the LEO landscape. Due to private companies’ goals to meet consumer demand, there is sometimes the opportunity cost of the development of COLA architecture. Although there has been no major orbital collision so far in 2026 that has been publicly attributed to failure of conjunction assessment, the operational burden associated with preventing one is rapidly increasing. In a 6-month period (December 2025 – May 2026), Starlink alone conducted more than 200,000 COLA manoeuvres [4].
The central research question is therefore investigating whether a supervised ML model can act as an effective first stage filter for conjunction screening. Rather than replace high-fidelity propagation or the subsequent calculation of collision probability, the proposed model is designed to identify satellite pairs that warrant further analysis at a much earlier stage than is currently industry standard of 7 days for LEO orbits [5]. Pairs classified as significant can then be passed to more computationally expensive screening processes, while those identified as low risk can be deprioritised. This transforms the problem, then. With limited computational resources, the adoption of ML to identify patterns in historical conjunction data significantly reduces the computational workload. This paper studies the trade-offs between accuracy and time in the COLA framework of the NewSpace era, as discussed in ref. [3].
To test this approach, two-line element (TLE) data is taken from historical conjunction data messages (CDMs) and used to construct examples of genuine conjunction events. These are compared against strictly selected negative pairs, representing objects that occupy broadly comparable orbital regions but remain sufficiently separated when propagated. The ML models then use the orbital relationships alone to predict whether the pairs should recall a CDM or not. These models are then evaluated, not only by their overall predictive performance, but also their ability to retain dangerous encounters. This distinction is essential because a missed conjunction is considerably more consequential than an unnecessary secondary evaluation. 
Returning to the central research question, it is hypothesised that the relationships contained within these orbital parameters are sufficiently distinct for an ML model to successfully reduce the number of candidate pairs whilst maintaining a high recall of genuine conjunction events. If supported, this would demonstrate the potential for ML to operate within existing space situational awareness (SSA) architectures, complementing established COLA methods.


\section{Current state of the art}\label{s:section}

LEO has experienced a vast expansion in the number of tracked space objects occupying its space in recent decades. An obvious reason for this is that space is becoming increasingly accessible. In 2026, space is being utilised as a tool for enterprise and space technology is seen in almost every sector; maps, Wi-Fi, and even healthcare all use technology that has been developed for space [6]. Despite this, a significant portion of the increased number of tracked objects in LEO comes from the detection of smaller objects that were previously undetected and therefore untracked. These objects are typically classified as orbital debris, which can include fragments of spacecraft, and inactive satellites, and accounts for 83\% of all tracked space objects [7]. Space exploration has not peaked yet either. It is predicted that 20,000 more satellites will be launched into LEO in the next decade [3]. 
Due to this growth, our orbits have become heavily congested and are constantly getting busier, causing satellite owners/operators (O/Os) to pour millions of pounds into COLA to protect increasingly valuable orbital infrastructure from a mission-ending collision. In the SSA industry, a number of processes are currently in place to protect the satellites in orbit from potential conjunctions. As a result, considerable research has focussed on improving the safety. Manfletti et al. describes the current situation in terms of its financial cost to O/Os—up to \$15M in a single year to protect a 300-satellite constellation (which would be on the smaller end of constellations) [8]. However, even this cost is conservative. 
For continuation…


\section{Materials and methods}\label{s:methods}

To investigate the effectiveness of an ML model as a first-stage filter for conjunction screening, the primary objective was to find how reliably the model could identify satellite pairs associated with genuine CDM generation using orbital data alone. A particular focus was on the model’s recall rate; a high recall rate means that the model could successfully identify the satellites that should generate a CDM to a high degree of accuracy, therefore minimising the number of potentially dangerous events incorrectly filtered out. 
Historical CDM data was used to provide real-world examples of satellite pairs that had previously flagged a conjunction assessment. Thousands of genuine CDMs were parsed to find pairs of satellites that we know generated a CDM using the current industry method [9]. Historical TLE data from both tracked space objects, which dates back to immediately before the CDM was generated, was then compiled into a positive dataset. In addition to this, a corresponding negative dataset was generated from further TLE data by selecting satellite pairs that satisfied defined non-conjunction criteria and did not produce a CDM under those conditions.
The CDM/no-CDM labels were withheld, leaving only orbital characteristics as the features used by the model in testing. After training the model, a fresh positive and negative testing dataset was then passed through the model. Its predictions were then classified according to whether it was likely to result in CDM generation. These were then compared with the known historical labels, allowing recall and other metrics to be calculated. The model was then evaluated on these metrics to determine its suitability as an initial conjunction screening filter. Outlined in this section are the stages of developing and testing this model.

\subsection{Data Retrieval, Processing, and Labelling}

A large component of the initial programme setup was all about collecting and transforming data. The key data required at this stage were historical CDMs and their time-appropriate TLE data. The necessity for time-appropriate TLE data is crucial. TLE data from after the conjunction assessment may include later observations or manoeuvres that were unavailable at the intended prediction time, potentially producing an unrealistic estimate of the model’s predictive performance.
To retrieve the necessary data (historical CDMs and historical TLEs), an automated session request is made through the Space-Track API [spacetrack website reference] and the specific data is downloaded as a JSON file. This is converted into a CSV file for use with the Pandas data analysis package. Historical data were retrieved in smaller batches to reduce request sizes and avoid timeouts. In addition to this, pauses were added between requests to respect the request-rate limit put in place by Space-Track. Successfully downloaded batches were then saved to a local manifest file, allowing the retrieval to continue after an interruption without repeating completed packets. 
Once all JSON files have been imported and converted appropriately, using the python data analysis library Pandas, the retrieved data was then processed into clean CSV files ready for labelling. In order to produce the positive testing dataset, the retrieved data was converted into one file which contained pairs of satellites with both the CDM data and the first TLE data from immediately before the creation time of the CDM for both satellites. This dataset is all labelled with a ‘1’ indicating that these satellite pairs have resulted in the generation of a CDM. 
Each CDM has a pair key, and this is used to prevent the combination AB being repeated under a new BA check, using up unnecessary computation. These pair keys allow later programmes to directly find key information about specific satellites, including all TLE-present orbital features. A separate programme compiled these features into a single packet associated with each pair, ready to directly pass into an ML model. These packets contained 27 distinct orbital features for every single tracked object. 
Finally, a general TLE pull request was made. A file containing several thousand individual tracked objects and their orbital data was produced. This was processed in a particular way, such that a search of the positive test dataset was made to ensure that the only TLEs left behind were those from tracked objects that have not generated a CDM in the past 3 months. This is the first step to generate our negative test dataset. Each satellite was randomly paired with another from that set. Just because these objects haven’t generated a CDM yet, does not mean that (when paired alongside a second arbitrary object) they wouldn’t produce one in the near future. As a result, these pairs cannot yet be labelled with a ‘0’. This label suggests that no CDM would be generated, which may not be true.

\subsection{Propagation and Strict Negatives}

To avoid this labelling problem for negative pairs, further propagation is required across a 7-day period; this 7-day period is taken directly from industry standards [9].

\section{Section with math}\label{s:math}

This section displays a number of mathematical expressions to showcase the math fonts used in the template.

\subsection{Roman letters} 

The Roman characters in math are just the same as the characters in the text---but in \textit{italic}. Here are some small letters: 
\begin{equation*}
a\{p - r \times l\} + \frac{w(g/z+j)}{i(t)+j(t)+k(t) - e^p - x^j} + \frac{h[f]+x^f}{k[y]-e^y + x^y} = f(j)^{6+y} - i^3_{g,j,p} \approx (p_{ji})^5.
\end{equation*}
Here are some capital letters as well: 
\begin{equation*}
G[p + P^7-Q] - A_B + L\{j\} = F(X) \to [Y+K]\times Z_f/H - [g_4 - i],\;\text{for any}\;i.
\end{equation*}
The punctuation is also the same in math as in text. 

\subsection{Digits}

The digits in math are also taken from the text font, just like the Roman characters. So 1, 2, 3, 4, 89, 03 (in text mode) is just the same as $1,\,2,\,3,\,4,\,89,\,03$ (in math mode). The percentage sign is also the same so 27\% (in text mode) is the same as $27\%$ (in math mode).

\subsection{Greek letters}

Here are some small Greek letters in a math display: 
\begin{equation*}
\alpha^{\theta} + \gamma^4 + g(a) - \zeta \times \frac{y(\lambda)}{\kappa \cdot k+ \sigma - s} \to \nu_{\eta} = \beta^\epsilon \cdot \delta - \mu + \frac{\xi^4}{\zeta_{ij}} \to m(x) + b^e + \omega .
\end{equation*}
Here are some capital Greek letters in another math display: 
\begin{equation}
F(\Psi) - G(\Phi) \times \Delta^{10} = \Sigma \cdot \Omega^2 - \frac{\Lambda_i}{\Theta_j} \cdot \frac{\Gamma(x)}{\Pi(t)}.
\label{e:capitalgreek}\end{equation}

Some but not all the variants of the standard Greek letters are available. There is a variant for theta ($\vartheta \neq \theta$), for epsilon ($\varepsilon \neq \epsilon$), for pi ($\varpi \neq \pi$), and for phi ($\varphi \neq \phi$). But there is no available variant for sigma ($\varsigma = \sigma$), rho ($\varrho = \rho$), and kappa (\texttt{\textbackslash varkappa} gives an error).

\subsection{Calligraphic letters} 

Calligraphic letters are taken from the same family of fonts as the Greek letters. The calligraphic letters are only available for capital letters. Here are some of the calligraphic letters: 
\begin{equation*}
\frac{\mathcal{A}(A)}{\mathcal{H}} - \mathcal{B}(\beta) + \mathcal{Q}(j+3f) \to \mathcal{Y} + \mathcal{M}(\mathcal{N} - 2) = \mathcal{J}(X) \neq \mathcal{X}(J) - \mathcal{F}[\mathcal{D} \times \Delta]^8. 
\end{equation*}

\subsection{Blackboard-bold letters} 

Here are some blackboard-bold letters, used in some mathematical expressions: $x \in\mathbb{R}^n$ but $y\in \mathbb{Q}$ and $[A_1,A_2] \in \mathbb{Z} \times \mathbb{N}$; so that $\mathbb{E}[X\mid Y] = Y + 3 \cdot \mathbb{P}(X^2) \notin \mathbb{C} \times \mathbb{1}(X \subseteq \mathbb{B})$.

\subsection{Symbols} 

The template comes with a full set of mathematical symbols. Here is a random assortment of some symbols: 
\begin{equation*}
\left(\frac{\sum_i \vec{z_i}}{\sum_{j=-\infty}^{+\infty} X_j}\right) \pm Q \subset \mathbb{R} \neq \mathbb{C} \implies \sqrt{\frac{\mathcal{X}}{2}} > \sqrt{\mathcal{Y}} \gg \hat{b}(y_0).
\end{equation*}

Lorem ipsum dolor sit amet, consectetur adipiscing elit. Sed aliquam magna vel urna ultrices, sed lacinia nulla mattis. Fusce non nunc nec est mollis malesuada. Nullam dignissim nulla sit amet libero facilisis, eget fringilla libero sagittis. Suspendisse potenti. Vivamus fermentum consectetur ante, at rhoncus nisi tristique vel:
\begin{equation*}
\sum_i \left[\frac{\hat{x}(i)}{\mathcal{Y}(i)}\right]^2 \leq 2\pi \Rightarrow \prod_j \bar{y}(j)^2 \geq \left\{\iint \frac{g(h)}{\gamma(s) \times \Delta(h)}\,dh\,ds \right\} \mapsto f\circ g(x) \iff \tilde{Q}(x) \propto \frac{\dot{x}}{2}.
\end{equation*}

\subsection{Bold characters}

In the template it is possible to bold all math characters. Roman characters can be bolded, such as $\bm{a} + \bm{D} = \bm{E}^2 + \bm{j}/\bm{i}$. Greek characters can also be bolded, such as $\bm{\alpha} + \bm{\Delta} = \bm{\epsilon}^2 + \bm{\Lambda}/\bm{\Phi}$. It is also possible to bold digits: $1 + 2 \neq \bm{1} + \bm{2}$. Finally, it is possible to bold calligraphic letters: $\bm{\mathcal{C}} + \bm{\mathcal{E}} - [\bm{\mathcal{X}}+\bm{\mathcal{Y}}]$. Blackboard-bold letters are already ``bold'' so they cannot be bolded further.

\subsection{Operators} 

Operators in mathematical expressions are typeset with the text font. Here are a few examples: 
\begin{equation*}
\max_x \pi(x) = \frac{\int \ln(u)du}{\int\exp(v)\,dv} \mapsto \cos(\theta) + \sin(\omega) - \min_z[\tan(z)].
\end{equation*}
Pellentesque habitant morbi tristique senectus et netus et malesuada fames ac turpis egestas. Aliquam erat volutpat. Donec tincidunt, quam sed pellentesque fermentum, lorem dolor consequat lectus, at pulvinar arcu ipsum in ligula. Mauris pretium ipsum id sapien posuere vehicula. 

\subsection{Some equations} 

This subsection illustrates how the different math fonts fit together. The equations below combine various characters together.

First, here are equations in text. Lorem ipsum dolor sit amet, consectetur adipiscing elit: $\prod_i\alpha_i(X) \geq 2 \cdot [\partial a_u(Z)/\partial Z] \propto \mathcal{X} \sim 5.2\pi \leadsto +\infty$. Lorem ipsum dolor sit amet, consectetur adipiscing elit: $\sum_{x=1}^4 \xi^x(j)=0, \; \forall j\in \mathbb{R}$. Here are some other expressions: $\mathcal{R}(z) = \max\{T_1,\ldots,T_{n}\} > \hat{z}$, and  $\mathbb{P}(T_n < z) = 1 - S(z) = F(z)$.  Sed vehicula ultrices odio, vitae maximus justo convallis at. Fusce quis consequat arcu: $\alpha^* = S(z^*)$ and $S^*(z^*)=\mathbb{E}[\bm{\alpha}\mid \beta]$. Integer non ante eget purus aliquet consequat. Mauris ut dui nec ligula consectetur commodo.

Here are equations in display mode, without labels:
\begin{align*}
\frac{\mathcal{S}(u^*)}{1 - \frac{\partial \ln(\theta)}{\partial \theta} + \left[\gamma^\eta\right] - \mathcal{G}(z^*)} &= \sum_{i=0}^{+\infty}\frac{d F^{i}(x)}{d x} - \Lambda(i) \geq \left\lbrace\sum_{i=0}^{+\infty}G^{i} - \Omega(i)\right\rbrace\\
\frac{\mathbb{Q}(z^*)}{1 -\mu + \exp(\xi) \times \mathbb{M}(z^*)} &= \prod_{i=0}^{+\infty}F^{i} - \Phi(i) \neq \left(\int_{0}^{+\infty}\frac{G^{i}}{\epsilon(i)} di\right).
\end{align*}
Lorem ipsum dolor sit amet, consectetur adipiscing elit. Sed vehicula ultrices odio, vitae maximus justo convallis at. Fusce quis consequat arcu. Integer non ante eget purus aliquet consequat. Mauris ut dui nec ligula consectetur commodo.

And here is a single equation in display mode without a label:
\begin{equation*}
\dot{y}(t) \to \sum_{j=0}^{+\infty}\mathcal{W}^j(z^{\ast}(t)) = \frac{S(z) \mathbb{E}(N(z))}{\mu_1 - \zeta_2} \cdot 2 - \left[\prod_{i=0}^{+\infty}\mathcal{F}^{i}\right]-\iiint\exp(\lambda(t) \mu(s)) \sin(\theta)dt\,ds\,d\theta.
\end{equation*}
Pellentesque habitant morbi tristique senectus et netus et malesuada fames ac turpis egestas. Nam consequat, ipsum eget tincidunt aliquam, mi est pellentesque lacus, nec ullamcorper leo magna a dui.

\subsection{Theorems, propositions, lemmas, and other results}

Here is a proposition with some more math:   
\begin{proposition}
Lorem ipsum dolor sit amet, consectetur adipiscing elit:
\begin{equation}
\sum_k\bm{S}_{k_x}(z) \approx \frac{S(z)^x}{k / 23 -\zeta\gamma [45- S(z)] + \ln(y) - j^2+x(l)}.
\label{e:type1}\end{equation}
Ut enim ad minim veniam, sunt in culpa qui officia deserunt mollit anim id est laborum. Excepteur sint occaecat cupidatat non proident, sunt in culpa qui officia deserunt mollit anim id est laborum.
\label{p:type1}\end{proposition}
\begin{proof} 
Here is the proof to the proposition. Donec commodo justo a eros malesuada, eget vulputate tortor accumsan. Sed ac pulvinar nulla. Etiam quis felis dapibus, vulputate metus eu, finibus nunc. Sed vel sodales dui. Nam venenatis dolor non orci tempus fermentum. Vivamus sodales justo a ligula cursus aliquet. Sed fringilla nunc vitae justo finibus, id placerat lectus sodales.
\end{proof} 

Vestibulum ante ipsum primis in faucibus orci luctus et ultrices posuere cubilia curae; Sed ac eros vel felis vehicula vehicula. Phasellus interdum justo a felis congue, vel dapibus est tristique. Now here is a lemma---note that all types of results have their own counter:
\begin{lemma} 
Ut enim ad minim veniam, quis nostrud exercitation ullamco laboris nisi ut aliquip ex ea commodo consequat:
\begin{equation}
z^* = \int_{0}^{\infty} \alpha(i) \cdot \frac{1-\beta}{1-\alpha(i)\beta}\,di.
\label{e:cv}\end{equation}
\label{p:cv}\end{lemma}

Quia dolor sit amet consectetur adipiscing velit, sed quia non numquam eius modi tempora incidunt, ut labore et dolore magnam aliquam quaerat voluptatem. Then here is another lemma:
\begin{lemma}
Consectetur adipiscing elit, sed do eiusmod tempor incididunt ut labore et dolore magna aliqua $\mathcal{Z}(\alpha)$. 
\begin{equation*}
\frac{\sum_i z^i}{\prod_i q^i} \to \frac{\int_{0}^{\infty} \alpha(i) \cdot [1-\beta]\,di}{1-\exp(\alpha)\sin(\beta)}.
\end{equation*}
Duis aute irure dolor in reprehenderit in voluptate velit esse cillum dolore eu fugiat nulla pariatur. 
\end{lemma}

Pellentesque habitant morbi tristique senectus et netus et malesuada fames ac turpis egestas. Aliquam erat volutpat. Donec tincidunt, quam sed pellentesque fermentum, lorem dolor consequat lectus, at pulvinar arcu ipsum in ligula. Mauris pretium ipsum id sapien posuere vehicula. Vivamus consectetur, turpis vitae pellentesque molestie, libero odio venenatis mauris, id vestibulum nisl nulla sit amet enim. Cras ultricies consectetur ex, sit amet vestibulum lacus hendrerit vel. Sed lacinia eleifend purus. Nunc dignissim eros et sapien lacinia eleifend. Curabitur sit amet velit eu metus suscipit tempus. Maecenas in libero felis. Sed quis eros sapien. Suspendisse potenti. Curabitur porta est sed enim efficitur, a tempor justo fermentum. Vivamus pulvinar enim vitae semper rutrum. Aenean non mauris sed lectus vestibulum feugiat eget et ipsum. Donec semper vitae enim at ultricies. Sed eu arcu justo.

Nemo enim ipsam voluptatem, quia voluptas sit, aspernatur aut odit aut fugit, sed quia consequuntur magni dolores eos, qui ratione voluptatem sequi nesciunt, neque porro quisquam est, qui dolorem ipsum. And here is a theorem following the lemma:
\begin{theorem} 
Lorem ipsum dolor sit amet, consectetur adipiscing elit, sed do eiusmod tempor incididunt ut labore et dolore magna aliqua:
\begin{equation*}
\mathbb{E}(N(z^*)) \approx \frac{1-\mathbb{P}(\alpha\pi)}{1-\pi}- \frac{f(y)}{z(p)^*} + P(\Gamma).
\end{equation*}
Ut enim ad minim veniam, quis nostrud exercitation ullamco laboris nisi ut aliquip ex ea commodo consequat.
\end{theorem}

Nemo enim ipsam voluptatem, quia voluptas sit, aspernatur aut odit aut fugit, sed quia consequuntur magni dolores eos, qui ratione voluptatem sequi nesciunt, neque porro quisquam est, qui dolorem ipsum. And here is a corollary following the theorem:
\begin{corollary} 
Lorem ipsum dolor sit amet, consectetur adipiscing elit, sed do eiusmod tempor incididunt ut labore et dolore magna aliqua:
\begin{equation*}
\mathcal{Q}^* = \int_{-\infty}^{+\infty} \mu(j)^2 \cdot \frac{1-\ln(\nu(j))}{1-\exp(\beta(j))}\,dj.
\end{equation*}
Ut enim ad minim veniam, quis nostrud exercitation ullamco laboris nisi ut aliquip ex ea commodo consequat.
\end{corollary}

\section{Section with lists}\label{s:lists}

Lists are not very common in scientific papers, but they are sometimes useful. This section presents a few possible lists.

\subsection{Subsection with itemized lists}


Here is a basic itemized list with one level:
\begin{itemize}
\item Et harum quidem rerum facilis est et expedita distinctio.
\item Nam libero tempore, cum soluta nobis est eligendi optio.
\item Temporibus autem quibusdam et aut officiis debitis aut rerum necessitatibus saepe eveniet, ut et voluptates repudiandae sint et molestiae non recusandae.
\item Nam libero tempore, cum soluta nobis est eligendi optio.
\end{itemize}
Et harum quidem rerum facilis est et expedita distinctio. Nam libero tempore, cum soluta nobis est eligendi optio, cumque nihil impedit, quo minus id, quod maxime placeat, facere possimus, omnis voluptas assumenda est, repellendus. Here is a list with two levels:
\begin{itemize}
\item Et harum quidem rerum facilis est et expedita distinctio.
\item Nam libero tempore, assumenda est, repellendus, cum soluta nobis est eligendi optio:
\begin{itemize}
\item Cumque nihil impedit
\item Quo minus id, quod maxime placeat. Emporibus autem quibusdam et aut officiis debitis aut rerum necessitatibus saepe eveniet---nam libero tempore, cum (soluta) nobis est eligendi optio.
\item Facere possimus, omnis voluptas assumenda est
\end{itemize}
\item Emporibus autem quibusdam et aut officiis debitis aut rerum necessitatibus saepe eveniet. Nam libero tempore, cum soluta nobis est eligendi optio.
\end{itemize}
Lorem ipsum dolor sit amet, consectetur adipisicing elit, sed do eiusmod tempor incididunt ut labore et dolore magna aliqua. Ut enim ad minim veniam, quis nostrud exercitation ullamco laboris nisi ut aliquip ex ea commodo consequat. Duis aute irure dolor in reprehenderit in voluptate velit esse cillum dolore eu fugiat nulla pariatur. Excepteur sint occaecat cupidatat non proident, sunt in culpa qui officia deserunt mollit anim id est laborum.

\subsection{Subsection with numbered lists}\label{s:numberedlists}

And here is a numbered lists with two levels:
\begin{enumerate}
\item Et harum quidem rerum facilis est et expedita distinctio.
\item Nam libero tempore, cum soluta nobis est eligendi optio.
\item Emporibus autem quibusdam et aut officiis debitis aut rerum necessitatibus saepe eveniet. Nam libero tempore, cum soluta nobis est eligendi optio:
\begin{enumerate}
	\item Cumque nihil impedit.
	\item Quo minus id, quod maxime placeat.
	\item Facere possimus, omnis voluptas assumenda est.
\end{enumerate}
\item Et harum quidem rerum facilis est et expedita distinctio.
\item Nam libero tempore, cum soluta nobis est eligendi optio.
\end{enumerate}
Itaque earum rerum hic tenetur a sapiente delectus, ut aut reiciendis voluptatibus maiores alias consequatur aut perferendis doloribus asperiores repellat. 

\section{Section with graphs}\label{s:graphs}

Here is a section with a variety of graphs. At vero eos et accusamus et iusto odio dignissimos ducimus, qui blanditiis praesentium voluptatum deleniti atque corrupti, quos dolores et quas molestias excepturi sint, obcaecati cupiditate non provident, similique sunt in culpa, qui officia deserunt mollitia animi, id est laborum et dolorum fuga. 

\subsection{Subsection with graphs at the top of the page}

A simple two-panel graph is on figure \ref{f:graph1}. It will be placed at the top of the page, just about here. Et harum quidem rerum facilis est et expedita distinctio. Nam libero tempore, cum soluta nobis est eligendi optio, cumque nihil impedit, quo minus id, quod maxime placeat, facere possimus, omnis voluptas assumenda est, omnis dolor repellendus. Temporibus autem quibusdam et aut officiis debitis aut rerum necessitatibus saepe eveniet, ut et voluptates repudiandae sint et molestiae non recusandae.

\begin{figure}[t]
\subcaptionbox{A first panel, 1951--2019\label{f:panelA}}{\includegraphics[scale=0.2,page=1]{\pdf}}\hfill
\subcaptionbox{A second panel, 1951--2019\label{f:panelB}}{\includegraphics[scale=0.2,page=2]{\pdf}}
\caption{Graph with two panels}
\note{This is a note for the graph. Nam libero tempore, cum soluta nobis est eligendi optio, cumque nihil impedit, quo minus id, quod maxime placeat, facere possimus.}
\label{f:graph1}\end{figure}

\subsection{References to figures and panels} 

As usual with LaTeX, it is easy to refer to a figure: see figure \ref{f:graph1}. It is possible to refer to a specific panel in a figure, for instance figure \ref{f:panelA} or figure \ref{f:panelB}. Its also possible to refer to the entire figure, for instance figure \ref{f:graph1} or figure \ref{f:graph2}.

\subsection{A subsection with a full-page figure}

Here is a full-page figure (figure \ref{f:graph2}). It will be placed in a full page about here. At vero eos et accusamus et iusto odio dignissimos ducimus, qui blanditiis praesentium voluptatum deleniti atque corrupti, quos dolores et quas molestias excepturi sint, obcaecati cupiditate non provident, similique sunt in culpa, qui officia deserunt mollitia animi, id est laborum et dolorum fuga. Et harum quidem rerum facilis est et expedita distinctio.

\begin{figure}[p]
\subcaptionbox{A first black-and-white panel}{\includegraphics[scale=0.2,page=5]{\pdf}}\hfill
\subcaptionbox{A second black-and-white panel}{\includegraphics[scale=0.2,page=6]{\pdf}}\\
\subcaptionbox{A third black-and-white panel}{\includegraphics[scale=0.2,page=7]{\pdf}}\hfill
\subcaptionbox{A fourth black-and-white panel}{\includegraphics[scale=0.2,page=8]{\pdf}}\\
\subcaptionbox{A fith black-and-white panel}{\includegraphics[scale=0.2,page=9]{\pdf}}\hfill
\subcaptionbox{A sixth black-and-white panel}{\includegraphics[scale=0.2,page=10]{\pdf}}
\caption{Graph with six black-and-white panels}
\note[Note]{Nam libero tempore, cum soluta nobis est eligendi optio, cumque nihil impedit, quo minus id, quod maxime placeat, facere possimus, omnis voluptas assumenda est, omnis dolor repellendus. Aenean fermentum purus id lacus volutpat, a eleifend mi posuere! Mauris nec nunc commodo, vehicula enim nec; vehicula ex. Nullam euismod lorem at eros efficitur: ut ultricies ante fringilla? Nam sagittis sapien id tortor commodo—a pulvinar velit ultricies. Integer ac magna velorci mollis vestibulum. Fusce id ipsum vel magna placerat vehicula. Curabitur ac lobortis justo.}
\label{f:graph2}\end{figure}

Nam libero tempore, cum soluta nobis est eligendi optio, cumque nihil impedit, quo minus id, quod maxime placeat, facere possimus, omnis voluptas assumenda est, omnis dolor repellendus. Temporibus autem quibusdam et aut officiis debitis aut rerum necessitatibus saepe eveniet, ut et voluptates repudiandae sint et molestiae non recusandae. Itaque earum rerum hic tenetur a sapiente delectus, ut aut reiciendis voluptatibus maiores alias consequatur aut perferendis doloribus asperiores repellat. 

\section{A section with table}

Here is a section with a variety of tables. Temporibus autem quibusdam et aut officiis debitis aut rerum necessitatibus saepe eveniet, ut et voluptates repudiandae sint et molestiae non recusandae. Itaque earum rerum hic tenetur a sapiente delectus, ut aut reiciendis voluptatibus maiores alias consequatur aut perferendis doloribus asperiores repellat. 

\subsection{A subsection with a simple table}

Table \ref{t:table1} is a simple table with one panel. Temporibus autem quibusdam et aut officiis debitis aut rerum necessitatibus saepe eveniet, ut et voluptates repudiandae sint et molestiae non recusandae. Itaque earum rerum hic tenetur a sapiente delectus, ut aut reiciendis voluptatibus maiores alias consequatur aut perferendis doloribus asperiores repellat. Itaque earum rerum hic tenetur a sapiente delectus, ut aut reiciendis voluptatibus maiores alias consequatur aut perferendis doloribus asperiores repellat. 

\begin{table}[t]
\caption{Basic table with one panel and multicolumns}
\begin{tabular*}{\textwidth}[]{p{4cm}@{\extracolsep\fill}cccc}
\toprule
& \multicolumn{2}{c}{Columns 1--2} & \multicolumn{2}{c}{Columns 3--4}\\
\cmidrule{2-3}\cmidrule{4-5}
& Column 1 &  Column 2 &  Column 3  &  Column 4 \\
\midrule
Line 1: & $\alpha$& $\beta$& $\gamma$ & $\delta$\\
Line 2: & $\epsilon$& $\phi$ & $\kappa$ & $\eta$  \\
Line 3: & $\kappa$& $\nu$& $\pi$ & $\kappa$ \\
Line 4: & $\psi$& $\mu$& $\nu$ & $\zeta$ \\
\bottomrule
\end{tabular*}
\note{This is a note for the table. Temporibus autem quibusdam et aut officiis debitis aut rerum necessitatibus saepe eveniet, ut et voluptates repudiandae sint et molestiae non recusandae. Pellentesque nec justo aliquet, commodo nulla sed, fringilla odio. Nullam non hendrerit nisi. Curabitur et metus vel velit blandit pharetra. Morbi interdum metus a erat bibendum, nec hendrerit eros ultricies. Vestibulum vel arcu id nulla ultricies commodo. Suspendisse potenti.}
\label{t:table1}\end{table}

Temporibus autem quibusdam et aut officiis debitis aut rerum necessitatibus saepe eveniet, ut et voluptates repudiandae sint et molestiae non recusandae. Itaque earum rerum hic tenetur a sapiente delectus, ut aut reiciendis voluptatibus maiores alias consequatur aut perferendis doloribus asperiores repellat.

\subsection{A subsection with a multi-panel table}

Table \ref{t:table2} is a more sophisticated table with several panels. Temporibus autem quibusdam et aut officiis debitis aut rerum necessitatibus saepe eveniet, ut et voluptates repudiandae sint et molestiae non recusandae. Itaque earum rerum hic tenetur a sapiente delectus, ut aut reiciendis voluptatibus maiores alias consequatur aut perferendis doloribus asperiores repellat. 

Suspendisse potenti. Ut venenatis maximus tellus, sit amet mattis lorem finibus nec. Mauris non libero eget ipsum ultricies congue id et eros. Duis nec vehicula nunc. In hac habitasse platea dictumst. Morbi non magna vitae ex fermentum placerat. Quisque eu ultrices velit. Nullam ac odio ac ex tincidunt posuere. Mauris viverra arcu eu metus ultrices, id vestibulum felis posuere.

Fusce eu magna in quam tincidunt placerat. Proin laoreet lacus eget lacinia ullamcorper. Nunc pulvinar, risus quis tempor vestibulum, odio odio fermentum tortor, nec dictum felis libero quis mauris. Vestibulum vel justo a eros efficitur fringilla. Sed at tempus urna. Nulla facilisi. Vivamus sit amet nisl at libero fermentum posuere.


\begin{table}[t]
\caption{Bigger table with several panels}
\begin{tabular*}{\textwidth}[]{p{2cm}@{\extracolsep\fill}ccccccc}
\toprule
    & Column 1 &  Column 2 &  Column 3  &  Column 4 &  Column 5 &  Column 6 &  Column 7 \\
\midrule
\multicolumn{8}{l}{A. A first panel with calligraphic letters}\\
Line 1: & $\mathcal{A}$ & $\mathcal{C}$ & $\mathcal{V}$  & -- & -- & $\mathcal{K}$ & $\mathcal{A}$\\
Line 2: & $\mathcal{X}$ &  $\mathcal{H}$ & $\mathcal{O}$  & -- & -- & $\mathcal{K}$ & $\mathcal{A}$  \\
\midrule
\multicolumn{8}{l}{B. A second panel with capital letters}\\
Line 3: & U & B & J  & K & A & K & A\\
Line 4: & N & Y & T  & L & T & K & A\\
Line 5: & G & S & Q  & P & Q & K & A\\
\midrule
\multicolumn{8}{l}{C. A third panel with blackboard-bold letters}\\
Line 6: & $\mathbb{U}$ & $\mathbb{B}$ & $\mathbb{J}$  & $\mathbb{K}$ & $\mathbb{K}$ & $\mathbb{K}$ & $\mathbb{A}$\\
Line 7: & $\mathbb{N}$ & $\mathbb{Y}$ & $\mathbb{T}$  & $\mathbb{L}$ & $\mathbb{T}$ & $\mathbb{U}$ & $\mathbb{E}$\\
Line 8: & $\mathbb{G}$ & $\mathbb{S}$ & $\mathbb{Q}$  & $\mathbb{P}$ & $\mathbb{Q}$ & $\mathbb{K}$ & $\mathbb{P}$\\
\midrule
\multicolumn{8}{l}{D. A fourth panel with numbers}\\
Line 9: & 1.0\% & 2.3\% & 4.5\%  & 9.0\% & 9.8\% & 91.2\% & 0\\
Line 10: & $\infty$ & 23 & 1 & 90 & 33 & 4.0\% & 0\\
Line 11: & -- & -- & 1.2  & 4.4 & 0 & 9.0\% & 0\\
\bottomrule
\end{tabular*}
\note[Note]{Temporibus autem quibusdam et aut officiis debitis aut rerum necessitatibus saepe eveniet, ut et voluptates repudiandae sint et molestiae non recusandae. Ut aut reiciendis voluptatibus maiores alias consequatur aut perferendis doloribus asperiores repellat.}
\note[Source]{\citet[table~1]{MS21} and \citet[table~2]{EMM21}.}
\label{t:table2}\end{table}

Temporibus autem quibusdam et aut officiis debitis aut rerum necessitatibus saepe eveniet, ut et voluptates repudiandae sint et molestiae non recusandae.

\section{A section with cross-references}

As usual, it is possible to reference an equation, such as equation \eqref{e:cv}. It is also possible to reference a section, such as section \ref{s:graphs}, or a subsection, such as section \ref{s:numberedlists}. It is also possible to reference an appendix, such as appendix \ref{a:appendix1}, or a subsection in an appendix, such as appendix~\ref{a:subappendix}. 

Of course it is possible to reference figures, such as figure \ref{f:graph2}, or tables, such as table~\ref{t:table1}. It is possible to refer to a specific panel in a figure, for instance figure \ref{f:panelA} or figure \ref{f:panelB}, and to the entire figure, for instance figure \ref{f:graph1} or figure \ref{f:graph2}.

\section{Quotes}

It is easy to quote material within the paper. Vivamus at commodo nisl. Integer sodales eros a metus efficitur, eu tincidunt leo ullamcorper. Sed fermentum tellus nec nisi lobortis, vel pharetra neque mollis. Cras non dolor ut ex semper lobortis. Vivamus sit amet sapien in elit sodales luctus. Donec a tortor ut elit consequat mollis ac id metus. Maecenas vehicula, turpis at lacinia interdum, dolor dolor pretium turpis, non fermentum purus magna at est. Quotes appear as follows:

\begin{quote}
Lorem ipsum dolor sit amet, consectetur adipisicing elit, sed do eiusmod
tempor incididunt ut labore et dolore magna aliqua. Ut enim ad minim veniam,
quis nostrud exercitation ullamco laboris nisi ut aliquip ex ea commodo
consequat. Duis aute irure dolor in reprehenderit in voluptate velit esse
cillum dolore eu fugiat nulla pariatur. Excepteur sint occaecat cupidatat non
proident, sunt in culpa qui officia deserunt mollit anim id est laborum.
\end{quote}

So this was a one-paragraph quote. Cras non dolor ut ex semper lobortis. Vivamus sit amet sapien in elit sodales luctus. Donec a tortor ut elit consequat mollis ac id metus. Maecenas vehicula, turpis at lacinia interdum, dolor dolor pretium turpis, non fermentum purus magna at est. 

\section{Conclusion}\label{s:ccl}

\paragraph{Summary}  At vero eos et accusamus et iusto odio dignissimos ducimus, qui blanditiis praesentium voluptatum deleniti atque corrupti, quos dolores et quas molestias excepturi sint, obcaecati cupiditate non provident, similique sunt in culpa, qui officia deserunt mollitia animi, id est laborum et dolorum fuga. At vero eos et accusamus et iusto odio dignissimos ducimus, qui blanditiis praesentium voluptatum deleniti atque corrupti, quos dolores et quas molestias excepturi sint.

\paragraph{Implications} Dolor sit amet, consectetur adipiscing elit. Nulla facilisi. Curabitur suscipit metus eget quam consequat, eget condimentum justo consectetur. Donec interdum ante nec felis fringilla, a rhoncus magna dapibus. Sed ultricies odio et magna consequat, ut posuere felis vehicula. Vivamus at commodo nisl. Integer sodales eros a metus efficitur, eu tincidunt leo ullamcorper. Sed fermentum tellus nec nisi lobortis, vel pharetra neque mollis. Cras non dolor ut ex semper lobortis. Vivamus sit amet sapien in elit sodales luctus. Donec a tortor ut elit consequat mollis ac id metus. Maecenas vehicula, turpis at lacinia interdum, dolor dolor pretium turpis, non fermentum purus magna at est. 

\paragraph{Next steps} Sed in purus nec nulla vulputate scelerisque. Curabitur rutrum aliquet sollicitudin. Suspendisse dapibus metus nunc, id tempus orci accumsan at. Sed rutrum purus velit, vel sollicitudin risus bibendum sed. Sed suscipit arcu ut purus malesuada dictum. Vivamus nec posuere neque. Sed consequat, odio sit amet cursus tincidunt, nulla nunc vulputate elit, ac pharetra eros metus nec ex. Vivamus at ligula id odio auctor blandit a et velit. Suspendisse potenti. Nunc efficitur est id tortor consectetur, sed convallis velit vulputate. Phasellus lacinia magna nec neque viverra ullamcorper.

\bibliography{\bib}
\appendix

% Enter appendix text:
\section{Generic appendix}\label{a:appendix1}

This is a generic appendix. Obcaecati cupiditate non provident, similique sunt in culpa, qui officia deserunt mollitia animi, id est laborum et dolorum fuga. Duis aute irure dolor in reprehenderit in voluptate velit esse cillum dolore eu fugiat nulla pariatur. Excepteur sint occaecat cupidatat non proident, sunt in culpa qui officia deserunt mollit anim id est laborum. 

\subsection{Generic subsection} 

This is a generic subsection in the appendix. Lorem ipsum dolor sit amet, consectetur adipisicing elit, sed do eiusmod tempor incididunt ut labore et dolore magna aliqua. Ut enim ad minim veniam, quis nostrud exercitation ullamco laboris nisi ut aliquip ex ea commodo consequat. 

In efficitur mi sed eros tincidunt tempus. Cras ut neque vehicula, convallis odio in, lacinia orci. Vivamus et nunc vestibulum, efficitur eros vitae, bibendum justo. Aenean tincidunt ligula ut augue pulvinar, ut placerat risus accumsan. Integer pellentesque odio ac commodo dictum. Ut ut justo a mi consequat vestibulum. 

Lorem ipsum dolor sit amet, consectetur adipiscing elit. Fusce gravida aliquam velit, sed fermentum odio fermentum eu. Donec eget fermentum ligula. Nulla facilisi. Integer vitae metus eget risus vulputate consectetur. Suspendisse potenti. Quisque dictum purus non justo aliquet, sed tristique odio dapibus. Sed tincidunt justo vitae leo mattis, a ullamcorper metus vestibulum. Sed ac sapien metus. Sed non mauris nec nisl vehicula iaculis eget eget enim. Suspendisse potenti. Sed ac purus volutpat, venenatis sem eget, fermentum lorem. Suspendisse eu convallis tortor. Suspendisse potenti. Nunc vel dictum libero. Cras eleifend, orci non faucibus gravida, sapien quam pharetra ligula, nec tincidunt elit ex at sapien. 

\subsection{Subsection with assumptions and results}

This subsection displays an assumption and results in the appendix to illustrate how they are typeset and numbered. First is an assumption:
\begin{assumption} 
Similique sunt in culpa, qui officia deserunt mollitia animi, id est laborum et dolorum fuga:
\begin{equation*}
\mathbb{E}(\Omega_{m,n}) = \mathbb{P}(\omega\cdot \mu - \xi) - \sum_{i=0}^{m}\sum_{j=1}^{n} \sigma(i,j) + 123^{56}.
\end{equation*}
Curabitur volutpat ultrices nunc id efficitur. Donec posuere, dui vel gravida elementum, ligula nisi lobortis elit, eu aliquet quam ligula id libero. Sed ac efficitur tur.
\label{a:appe}\end{assumption}

Lorem ipsum dolor sit amet, consectetur adipisicing elit, sed do eiusmod
tempor incididunt ut labore et dolore magna aliqua. Ut enim ad minim veniam,
quis nostrud exercitation ullamco laboris nisi ut aliquip ex ea commodo
consequat. Next is a theorem:
\begin{theorem}
Duis aute irure dolor in reprehenderit in voluptate velit esse
cillum dolore eu fugiat nulla pariatur. Excepteur sint occaecat cupidatat non
proident, sunt in culpa qui officia deserunt mollit anim id est laborum:
\begin{equation}
X^* =\iint_{0}^{\infty} \alpha(i) \cdot \mathcal{A}^2(i) + \mathbb{P}(X\mid Z(i))\,di\,dj,\end{equation}
integer vestibulum sapien nec velit varius, nec scelerisque nunc eleifend. Cras sollicitudin, justo sit amet tempus vehicula, ex ex vulputate nulla, non vestibulum quam tortor nec nisl. 
\end{theorem}

Lorem ipsum dolor sit amet, consectetur adipiscing elit. Sed aliquam magna vel urna ultrices, sed lacinia nulla mattis. Fusce non nunc nec est mollis malesuada. Here is another theorem based on assumption~\ref{a:appe}:
\begin{theorem}
Consectetur adipiscing elit, sed do eiusmod tempor incididunt ut labore et dolore magna aliqua: $y(\gamma) \geq 3\pi + \cos(\vartheta)$.
\label{t:theorem1}\end{theorem}

Lorem ipsum dolor sit amet, consectetur adipiscing elit. Sed aliquam magna vel urna ultrices, sed lacinia nulla mattis. Fusce non nunc nec est mollis malesuada. Nullam dignissim nulla sit amet libero facilisis, eget fringilla libero sagittis. Suspendisse potenti. Vivamus fermentum consectetur ante, at rhoncus nisi tristique vel. Vivamus in est quis justo fermentum lacinia ac eu leo. Maecenas nec tempor nisi---as in theorem \ref{t:theorem1}.

\subsection{A subsection with math}

Here is math and equations in the appendix---see equations \eqref{e:appendix1} and \eqref{e:experiments}. Temporibus autem quibusdam $\xi$ et aut officiis debitis aut rerum necessitatibus saepe eveniet ut et voluptates repudiandae sint et molestiae non recusandae $1-\gamma$. Itaque earum rerum hic $\mathcal{S}(\zeta^0)$ tenetur a sapiente delectus $\mathcal{B}^\theta$, ut aut reiciendis voluptatibus maiores alias consequatur aut perferendis doloribus asperiores repellat $\mathbb{V}^i$:
\begin{equation}
\mathbb{V}^r = (1-\gamma) \times 0 +\gamma S(z^*) v^s+\gamma [1-S(z^*)] \mathcal{V}^i-c.
\label{e:appendix1}\end{equation}

Ut enim ad minima veniam, quis nostrum exercitationem ullam corporis suscipit laboriosam, nisi ut aliquid ex ea commodi consequatur? Quis autem vel eum iure reprehenderit qui in ea voluptate velit esse quam nihil molestiae consequatur, vel illum qui dolorem eum fugiat quo voluptas nulla pariatur? 
\begin{equation}
\mathbb{E}(N(z)) = \frac{1}{1-\gamma F(z)}.
\label{e:experiments}\end{equation}

\section{Another appendix}\label{a:appendix2}

Here is a second appendix. At vero eos et accusamus et iusto odio dignissimos ducimus, qui blanditiis praesentium voluptatum deleniti atque corrupti.

\subsection{Larger figure, without panel, in the appendix} 

Here is a large, simple figure in the appendix (see figure \ref{f:appendix1}). At vero eos et accusamus et iusto odio dignissimos ducimus, qui blanditiis praesentium voluptatum deleniti atque corrupti, quos dolores et quas molestias excepturi sint, obcaecati cupiditate non provident, similique sunt in culpa, qui officia deserunt mollitia animi, id est laborum et dolorum fuga.

\begin{figure}[t]
\includegraphics[scale=0.3,page=8]{\pdf}
\caption{A larger graph in the appendix}
\note[Source]{This is the source for the larger graph. Nam libero tempore, cum soluta nobis est eligendi optio, cumque nihil impedit, quo minus id, quod maxime placeat, facere possimus.}
\label{f:appendix1}\end{figure}

Lorem ipsum dolor sit amet, consectetur adipisicing elit, sed do eiusmod
tempor incididunt ut labore et dolore magna aliqua. Ut enim ad minim veniam,
quis nostrud exercitation ullamco laboris nisi ut aliquip ex ea commodo
consequat. Duis aute irure dolor in reprehenderit in voluptate velit esse
cillum dolore eu fugiat nulla pariatur. Excepteur sint occaecat cupidatat non
proident, sunt in culpa qui officia deserunt mollit anim id est laborum.

Nullam fringilla, risus sit amet tincidunt sagittis, ligula nisi suscipit nisi, nec vulputate libero odio sit amet metus. Sed auctor elit nec orci venenatis, in dignissim lacus venenatis. Curabitur auctor odio sit amet leo molestie, vel lacinia sem congue. Integer id interdum metus. Ut ultrices ultricies lorem eget egestas. Integer efficitur libero id sapien posuere, at laoreet mauris consectetur. Nulla facilisi. Sed dapibus lectus at lacus interdum, ac vehicula dui bibendum. In hac habitasse platea dictumst. Sed luctus metus quis libero ultrices, vitae fringilla purus varius. Sed efficitur suscipit felis vel dapibus. Nunc condimentum laoreet ipsum, ut faucibus risus sodales sed. Sed dignissim, magna vel suscipit fermentum, lectus lorem suscipit tortor, eu suscipit elit tortor non dolor. 

\subsection{Subsection with references}\label{a:subappendix}

Here is a sentence with some references in text: \citet{MS19,MS21} found something but that is an uncommon result. The references can also go in parentheses at the end of the sentence \citep{MS21,MS22}. The references go to the reference list at the end of the main text---so appendix and main text share a common list of references.

Lorem ipsum dolor sit amet, consectetur adipisicing elit, sed do eiusmod
tempor incididunt ut labore et dolore magna aliqua. Ut enim ad minim veniam,
quis nostrud exercitation ullamco laboris nisi ut aliquip ex ea commodo
consequat. Duis aute irure dolor in reprehenderit in voluptate velit esse
cillum dolore eu fugiat nulla pariatur. Excepteur sint occaecat cupidatat non
proident, sunt in culpa qui officia deserunt mollit anim id est laborum.

\subsection{Subsection with a foonote}

Here is a sentence with a footnote.\footnote{Nemo enim ipsam voluptatem quia voluptas sit aspernatur aut odit aut fugit, sed quia consequuntur magni dolores eos qui ratione voluptatem sequi nesciunt.} The numbering of the footnotes, just like the numbering of the pages, continues from the main text to the appendix.

Lorem ipsum dolor sit amet, consectetur adipisicing elit, sed do eiusmod tempor incididunt ut labore et dolore magna aliqua. Ut enim ad minim veniam,
quis nostrud exercitation ullamco laboris nisi ut aliquip ex ea commodo
consequat. Duis aute irure dolor in reprehenderit in voluptate velit esse
cillum dolore eu fugiat nulla pariatur. Excepteur sint occaecat cupidatat non
proident, sunt in culpa qui officia deserunt mollit anim id est laborum.

\section{Last appendix with text}

This is a last appendix with a bit of text.

\subsection{A first subsection}

Sed vestibulum ex a tristique lacinia. Integer interdum magna vel magna rutrum fermentum. Maecenas sed mi in nunc convallis rutrum. Vivamus dapibus bibendum est, ac tincidunt ipsum fermentum id. Nunc id est turpis. Suspendisse potenti. Sed ac laoreet nulla, eu sollicitudin libero. Suspendisse tempus orci nec mauris volutpat eleifend. Integer nec tristique libero. Sed vehicula ipsum sit amet magna accumsan, nec pulvinar nisl consequat. In ac hendrerit turpis. Sed aliquet luctus mauris, vitae congue turpis. Nullam blandit lacus vel interdum eleifend. 

\subsection{A second subsection}

Nam pretium mauris eros, nec sollicitudin risus tincidunt a. Vivamus in augue vitae ligula scelerisque dapibus. Integer eget metus aliquet, efficitur nibh eget, pharetra eros. Suspendisse luctus interdum ex id suscipit. Donec quis augue mauris. Nulla ut erat eget nisl hendrerit malesuada. Vivamus eu tortor sit amet sem fringilla eleifend. In hac habitasse platea dictumst. Suspendisse potenti. Ut eget libero sed orci tempor ullamcorper:
\begin{equation}
S^*(z) = \frac{\alpha}{1-\gamma (1-\alpha)} > \int_{-\infty}^{+\infty}\alpha(x)\,dx.
\label{e:type1Classical}\end{equation}

\subsection{A third subsection}

Lorem ipsum dolor sit amet, consectetur adipiscing elit. Sed aliquam magna vel urna ultrices, sed lacinia nulla mattis. Fusce non nunc nec est mollis malesuada. Nullam dignissim nulla sit amet libero facilisis, eget fringilla libero sagittis. Suspendisse potenti. Vivamus fermentum consectetur ante, at rhoncus nisi tristique vel. Vivamus in est quis justo fermentum lacinia ac eu leo. Sed ac pulvinar nulla. Etiam quis felis dapibus, vulputate metus eu, finibus nunc. Sed vel sodales dui. Nam venenatis dolor non orci tempus fermentum. Vivamus sodales justo a ligula cursus aliquet. Sed fringilla nunc vitae justo finibus, id placerat lectus sodales.

\end{document}